% TOP_PROBLEM: {"id": 3079, "title": "Kneser–Poulsen conjecture", "category_id": 6, "category": {"id": 6, "name": "geometry", "display_name": "Geometry", "description": "Euclidean and non-Euclidean geometry, geometric structures.", "slug": "geometry", "order_index": 6, "created_at": "2026-07-31T15:26:25.671Z"}, "status": "open", "problem_number": "OPG-2089", "set_id": 12}
% FIRST_POSTED: 2026-10-01
% ATTEMPT_STATUS: unresolved
% COMPLETION: 5
% MODEL: GPT 6 Astra Ultra
% UnsolvedMath ID 3079; source catalogue code OPG-2089
% !TeX program = lualatex
% Research attempt, 2026-10-01. Canonical statement preserved verbatim.
\documentclass[11pt,a4paper]{article}
\usepackage[margin=25mm]{geometry}
\usepackage{fontspec}
\setmainfont{lmroman10-regular.otf}[BoldFont=lmroman10-bold.otf,
 ItalicFont=lmroman10-italic.otf,BoldItalicFont=lmroman10-bolditalic.otf]
\usepackage{amsmath,amssymb,mathtools}
\usepackage{unicode-math}
\setmathfont{latinmodern-math.otf}
\AtBeginDocument{\let\setminus\smallsetminus}
\usepackage{xurl,hyperref}
\hypersetup{colorlinks=true,urlcolor=blue}
\setcounter{secnumdepth}{0}
\setlength{\parindent}{0pt}
\setlength{\parskip}{5pt}
\setlength{\emergencystretch}{3em}
\begin{document}
\section*{Kneser–Poulsen conjecture}\label{3079}
{\small UnsolvedMath ID 3079 · OPG-2089 · Geometry · OpenGarden}
\subsection{Definitions and mathematical statement}
Let $n,N\ge1$ and consider two labelled configurations of centres
$p_1,\ldots,p_N$ and $q_1,\ldots,q_N$ in $\mathbb R^n$.
Repeated centres are allowed. For $x\in\mathbb R^n$, let
$B(x,1)=\{z:\|z-x\|_2\le1\}$ be the closed unit ball, and let
$\operatorname{vol}_n$ denote $n$-dimensional Lebesgue measure.
Assume the second configuration is an expansion of the first in the
pairwise-distance sense:
\[
 \|q_i-q_j\|_2\ge\|p_i-p_j\|_2
                    \qquad(1\le i,j\le N).
\]
The conjecture asserts
\[
 \operatorname{vol}_n\!\left(\bigcup_{i=1}^N B(q_i,1)\right)
 \ge
 \operatorname{vol}_n\!\left(\bigcup_{i=1}^N B(p_i,1)\right).
\]
The volume counts overlapping regions once, as ordinary volume of a
union, rather than as a sum of ball volumes. The same radius is assigned
to every ball both before and after the rearrangement; replacing one by
any other common positive radius gives an equivalent assertion by scaling.

Only the two endpoint configurations must satisfy the displayed distance
inequalities. The problem does not assume that they can be joined by a
continuous motion during which every pairwise distance is nondecreasing.
Such a hypothesis would restrict the class of rearrangements. Both
configurations must lie in the same specified Euclidean dimension.
The source entry asks about the union of equal balls; corresponding
questions for unequal radii or intersections are distinct extensions.
\subsection{Short English statement}
Can expanding all pairwise distances between centres ever decrease the volume covered by equal balls?
\subsection{Research attempt}
\paragraph{Scope and literature check.}
This attempt tests whether pairwise overlap accounting can prove the endpoint
expansion inequality. Bezdek--Connelly [A1] prove the planar theorem; their
result is prior work, not established anew here. The proof below handles an
arbitrary dimension when the initial positive-overlap graph is triangle-free.
No interpolation of the centres or differentiability of their motion is assumed.
All volumes below are Lebesgue volumes, so tangencies have zero volume.

\paragraph{Exact two-ball calculation.}
Write $\kappa_j$ for the volume of the unit ball in $\mathbb R^j$, with
$\kappa_0=1$. Two unit balls with centre distance $d\in[0,2]$ have intersection
volume
\[
 I_n(d)=2\kappa_{n-1}\int_{d/2}^1(1-t^2)^{(n-1)/2}\,dt.
\]
To obtain this, put their centres at $0$ and $d e_1$. Their intersection is
split by $x_1=d/2$ into two congruent spherical caps. The cap in the first
ball has sections of radius $\sqrt{1-t^2}$. Fubini's theorem gives the formula.
For $d\ge2$, put $I_n(d)=0$. For $0<d<2$ differentiation gives
\[
 I_n'(d)=-\kappa_{n-1}(1-d^2/4)^{(n-1)/2}\le0.
\]
Continuity handles $d=0,2$. This proves the conjecture for two balls in every
dimension, including coincident centres. For $n=1$ the formula is $I_1(d)=2-d$.

\paragraph{A triangle-free overlap theorem.}
Let $H_p$ have vertex set $[N]$ and edge $ij$ when $\|p_i-p_j\|<2$.
Suppose $H_p$ contains no triangle. No point belongs to the interiors of three
initial balls: such a point would force all three centre distances to be
strictly less than two by the triangle inequality. The union of the ball
boundaries has measure zero. Consequently, almost everywhere the multiplicity
$k$ of coverage is at most two, and
$\mathbf1_{k\ge1}=k-\binom{k}{2}$. Integration yields the exact identity
\[
 V(p)=N\kappa_n-\sum_{i<j}I_n(\|p_i-p_j\|).
\]
Every positive-overlap edge for the expanded configuration was already an
edge for the initial one. Thus $H_q\subseteq H_p$ is triangle-free as well,
and the same identity holds for $V(q)$. Monotonicity of $I_n$ now gives
$V(q)-V(p)\ge0$ term by term. This covers forests and also triangle-free graphs
with cycles; the acyclicity hypothesis sometimes proposed for this argument
is unnecessarily strong.

\paragraph{Stress test of the extension.}
For three balls with a common region of positive volume the correct formula
has an additional \emph{positive} triple-intersection term. Take three coincident
centres initially. Their union has volume $\kappa_n$, while the pairwise
truncation incorrectly returns zero. After moving them far enough apart,
the triple term vanishes. Thus triple-intersection terms can decrease under
expansion; dropping them is not a termwise proof of monotonicity. With more
balls the alternating higher intersections introduce the same obstruction.
The displayed pairwise method therefore proves a restricted theorem, not the
full conjecture even when every two-ball calculation is exact.

\paragraph{Verification and remaining gap.}
The proof explicitly checks coincident centres, tangencies, disappeared edges,
and the one-dimensional lens formula. Its only analytic inputs are Euclidean
slicing and finite additivity of integrals. No numerical volume experiment is
used as proof. A continuation must control higher intersections collectively,
or use an invariant that avoids inclusion--exclusion cancellation. Merely
knowing that individual pair overlaps shrink is insufficient.

\paragraph{Final status.}
\textbf{UNRESOLVED.} The triangle-free positive-overlap case is proved here;
no claim of novelty for this elementary special case is made. Arbitrary
initial overlap patterns in the catalogue's quantified dimensions remain
outside this attempt. Completion estimate: 5\%.

\subsection{Sources}
\begin{itemize}
\item[S1] ULAM.AI and the UnsolvedMath Contributors, supplied problems.json, record 3079 (OPG-2089), OpenGarden collection. Catalogue identity and source transcription; CC BY 4.0.
\url{https://huggingface.co/datasets/ulamai/UnsolvedMath}.
\item[S2] Open Problem Garden, Kneser–Poulsen conjecture. Original problem and discussion; the mathematical formulation below is expanded from the supplied source record.
\url{http://www.openproblemgarden.org/op/kneser_poulsen_conjecture}.
\item[A1] K. Bezdek and R. Connelly, \emph{Pushing disks apart---The Kneser--Poulsen conjecture in the plane}, arXiv:math/0108098. Primary abstract checked 1 October 2026. \url{https://arxiv.org/abs/math/0108098}.
\end{itemize}
\end{document}
